% eksperymentalne tłumaczenie na włoski

```latex
%

\section{Konstrukcje bez cyrkla} % lang-pl

\section{Costruzioni senza compasso} % lang-it

\begin{theorem}[twierdzenie Ponceleta-Steinera] % lang-pl
\begin{theorem}[teorema di Poncelet-Steiner] % lang-it
\label{thm:poncelet_steiner}%
\index{twierdzenie!Ponceleta-Steinera}%
    Jeśli dana konstrukcja jest wykonalna za pomocą cyrkla i linijki, to jest ona wykonalna za pomocą samej linijki, o ile dany jest na płaszczyźnie pewien okrąg wraz ze środkiem. % lang-pl

    Se una costruzione è realizzabile con riga e compasso, allora è realizzabile con la sola riga, purché nel piano sia dato un cerchio insieme al suo centro. % lang-it
    \index{twierdzenie!Ponceleta-Steinera}%
\end{theorem}

Nazwa twierdzenia pochodzi od dwóch panów: Jeana Ponceleta, który postawi w 1822 roku hipotezę oraz Jakoba Steinera, który dostarczy brakujący dowód w 1833 roku. % lang-pl

Il nome del teorema deriva da due studiosi: Jean Poncelet, che formulò la congettura nel 1822, e Jakob Steiner, che fornì la dimostrazione mancante nel 1833. % lang-it
\index[persons]{Poncelet, Jean}%
\index[persons]{Steiner, Jakob}%
Zamiast całego okręgu, wystarczy, że mamy do dyspozycji mały jego łuk -- to wzmocnienie, o które pokusi się Francesco Severi w 1904 roku. % lang-pl

Al posto dell'intero cerchio, è sufficiente avere a disposizione un suo piccolo arco: questo rafforzamento fu ottenuto da Francesco Severi nel 1904. % lang-it
\index[persons]{Severi, Francesco}%

Sama linijka to za mało, bo nie pozwala na wyciąganie pierwiastków kwadratowych, tak jak opisze to Łukasz Rajkowski w $\Delta_{15}^8$. % lang-pl

La sola riga non basta, poiché non permette di estrarre radici quadrate, come descrive Łukasz Rajkowski in $\Delta_{15}^8$. % lang-it
\index[persons]{Rajkowski, Łukasz}%
Ale zamiast środka okręgu możemy wymagać: % lang-pl

Tuttavia, al posto del centro del cerchio possiamo richiedere: % lang-it
\begin{itemize}
    \item równoległoboku, % lang-pl
    \item drugiego koncentrycznego okręgu, % lang-pl
    \item drugiego okręgu przecinającego pierwszy, % lang-pl
    \item drugiego okręgu rozłącznego z pierwszym i punktu na prostej łączącej ich środki albo ich osi potęgowej... % lang-pl
    \item un parallelogramma, % lang-it
    \item un secondo cerchio concentrico, % lang-it
    \item un secondo cerchio che interseca il primo, % lang-it
    \item un secondo cerchio disgiunto dal primo e un punto sulla retta che congiunge i loro centri oppure sul loro asse radicale... % lang-it
\end{itemize}

Zapewne istnieje jeszcze więcej takich warunków; o części napisze Eves \cite[s. 174-182]{eves1_1972}. % lang-pl

August Adler wspomniany we wcześniejszej sekcji zauważy, że linijka z dwoma brzegami albo ekierka pozwalają wykonać wszystkie klasyczne konstrukcje. % lang-pl
Esistono probabilmente molte altre condizioni di questo tipo; alcune sono descritte da Eves \cite[p. 174-182]{eves1_1972}. % lang-it
August Adler, menzionato nella sezione precedente, osservò che una riga a doppio bordo oppure una squadra consentono di eseguire tutte le costruzioni classiche. % lang-it


Dwie trójkątne linijki (w kształcie litery L) pozwalają podwoić sześcian (problem \ref{problem_delijski}) i dzielić kąt na trzy równe (problem \ref{trysekcja_kata}). % lang-pl
Due righe triangolari (a forma di L) permettono di duplicare il cubo (problema \ref{problem_delijski}) e di trisecare un angolo (problema \ref{trysekcja_kata}). % lang-it

\begin{problem}
    Dany jest okrąg $\Gamma$ ze środkiem $O$ oraz odcinek. % lang-pl
    
    Skonstruować środek odcinka. % lang-pl
    Sono dati un cerchio $\Gamma$ con centro $O$ e un segmento. % lang-it
    
    \hfill \emph{(15 kroków)}. % Hartshorne s. 192 % lang-pl
    Costruire il punto medio del segmento. % lang-it
    \hfill \emph{(15 passi)}. % Hartshorne s. 192 % lang-it
\end{problem}

\begin{problem}
    Dany jest okrąg $\Gamma$ ze środkiem $O$, prosta $l$ oraz punkt $P$. % lang-pl
    
    Skonstruować prostą równoległą do $l$, która przechodzi przez $P$. % lang-pl
    Sono dati un cerchio $\Gamma$ con centro $O$, una retta $l$ e un punto $P$. % lang-it
    
    \hfill \emph{(16 kroków)}. % Hartshorne s. 192 % lang-pl
    Costruire la retta parallela a $l$ passante per $P$. % lang-it
    \hfill \emph{(16 passi)}. % Hartshorne s. 192 % lang-it
\end{problem}

\begin{problem}
    Dany jest okrąg $\Gamma$ ze środkiem $O$, odcinek $OA$ oraz półprosta $OB$. % lang-pl
    
    Skonstruować punkt $C$ na półprostej taki, że $OA$ i $OC$ są przystające. % lang-pl
    Sono dati un cerchio $\Gamma$ con centro $O$, il segmento $OA$ e la semiretta $OB$. % lang-it
    
    \hfill \emph{(17 kroków)}. % Hartshorne s. 193 % lang-pl
    Costruire un punto $C$ sulla semiretta tale che $OA$ e $OC$ siano congruenti. % lang-it
    \hfill \emph{(17 passi)}. % Hartshorne s. 193 % lang-it
\end{problem}

\begin{problem}
    Dany jest okrąg $\Gamma$ ze środkiem $O$, prosta $l$ oraz punkt $P$. % lang-pl
    
    Skonstruować prostą prostopadłą do $l$, która przechodzi przez $P$. % lang-pl
    Sono dati un cerchio $\Gamma$ con centro $O$, una retta $l$ e un punto $P$. % lang-it
    
    \hfill \emph{(33 kroki)}. % Hartshorne s. 193 % lang-pl
    Costruire la retta perpendicolare a $l$ passante per $P$. % lang-it
    \hfill \emph{(33 passi)}. % Hartshorne s. 193 % lang-it
\end{problem}

Zetel \cite[s. 17]{zetel_2020} poda wariację tego problemu: jego prosta $l$ jest średnicą okręgu. % lang-pl

Premier Rosji Michaił Miszustin odwiedzi we wrześniu 2021 roku liceum imienia P. L. Kapicy w Dołgoprudnym i~poprosi uczniów o~rozwiązanie jeszcze prostszego zadania, gdzie punkt $P$ leży na okręgu $\Gamma$. % lang-pl
 % https://www.theguardian.com/science/2021/sep/20/did-you-solve-it-russias-prime-minister-sets-a-geometry-puzzle % lang-pl
Zetel \cite[p. 17]{zetel_2020} propone una variante di questo problema: la sua retta $l$ è un diametro del cerchio. % lang-it
Nel settembre 2021 il primo ministro russo Michail Mišustin visiterà il liceo P. L. Kapica di Dolgoprudnyj e chiederà agli studenti di risolvere un problema ancora più semplice, nel quale il punto $P$ giace sul cerchio $\Gamma$. % https://www.theguardian.com/science/2021/sep/20/did-you-solve-it-russias-prime-minister-sets-a-geometry-puzzle % lang-it

\begin{problem}
    Dany jest okrąg $\Gamma$ ze środkiem $O$, prosta $l$ oraz dwa punkty $A$, $B$. % lang-pl
    
    Skonstruować punkt, gdzie prosta $l$ przecina okrąg o środku $A$ i promieniu $AB$. % lang-pl
    Sono dati un cerchio $\Gamma$ con centro $O$, una retta $l$ e due punti $A$, $B$. % lang-it
    
    \hfill \emph{(54 kroki)}. % Hartshorne s. 193 % lang-pl
    Costruire il punto in cui la retta $l$ interseca il cerchio di centro $A$ e raggio $AB$. % lang-it
    \hfill \emph{(54 passi)}. % Hartshorne s. 193 % lang-it
\end{problem}

\begin{problem}
    Dany jest okrąg $\Gamma$ ze środkiem $O$ oraz punkt $A$ poza okręgiem. % lang-pl
    
    Skonstruować styczne do $\Gamma$ przechodzące przez $A$. % lang-pl
    Sono dati un cerchio $\Gamma$ con centro $O$ e un punto $A$ esterno al cerchio. % lang-it
    
    Costruire le tangenti a $\Gamma$ passanti per $A$. % lang-it
\end{problem}

\begin{problem}
    Dany jest okrąg $\Gamma$ ze środkiem $O$. % lang-pl
    
    Skonstruować pięciokąt foremny wpisany w $\Gamma$. % lang-pl
    È dato un cerchio $\Gamma$ con centro $O$. % lang-it
    
    \hfill \emph{(około 50 kroków)}. % Hartshorne s. 193 % lang-pl
    Costruire un pentagono regolare inscritto in $\Gamma$. % lang-it
    \hfill \emph{(circa 50 passi)}. % Hartshorne s. 193 % lang-it
\end{problem}

https://www.deltami.edu.pl/2017/06/z-sama-linijka-na-okrag-styczna-do-okregu-brakujacy-punkt/

https://www.deltami.edu.pl/2014/02/sama-linijka-mozna-nakreslic-okrag/

%
```
